\section{Related Work}
\label{sec:related}

\subsection{Off-Road Segmentation Datasets and Benchmarks}

\textbf{RUGD.} Wigness \etal{}~\cite{wigness2019rugd} released RUGD, more than
7,000 frames from video sequences recorded by a small unmanned ground robot
whose camera sits under 25\,cm above the ground, densely annotated with 24
classes. They partition the 18 sequences into train, validation, and test
splits, and deliberately place \emph{creek} (the only sequence with substantial
rock-bed terrain) and \emph{trail-7} (heavy motion blur) in the test split. Their
baselines all use a ResNet-50 encoder with PSPNet, UPerNet, or simpler
decoders, trained at $688\times555$. These reach 35.7--36.7\% mIoU on
validation but only 31.8--32.2\% on test (24-class mIoU; Tables~II and~III
of~\cite{wigness2019rugd}). Every baseline scores 0.0 IoU on \emph{bridge}
(Table~V of~\cite{wigness2019rugd}). We follow their sequence split exactly and
use their test mIoU as a historical reference point, not as a directly
comparable number, since our resolution, class averaging, and training budget
differ (Section~\ref{sec:methods}).

\textbf{RELLIS-3D.} Jiang \etal{}~\cite{jiang2021rellis} recorded RELLIS-3D with
a larger off-road vehicle on a Texas A\&M University campus. It provides
6,235 annotated camera images and 13,556 annotated LiDAR scans. On their image
test split, HRNetV2+OCR and Gated-SCNN reach only 51.55\% and 52.92\% mIoU
(Table~II of~\cite{jiang2021rellis}). The authors attribute the low scores to
severe class imbalance and unclear terrain boundaries. RELLIS-3D also contains
classes that RUGD lacks, such as \emph{mud}, \emph{puddle}, and \emph{rubble},
which motivates the shared group-level class map we use for the cross-dataset
test.

\textbf{Off-road methods.} GA-Nav~\cite{guan2022ganav} merges the fine classes
of both datasets into six navigability groups and adds a group-wise attention
mechanism to a transformer backbone. Under that grouping, GA-Nav reports
89.08\% mIoU on RUGD and 74.44\% on RELLIS-3D, compared with 55.88\% and
62.27\% for its DeepLabv3+ baseline (Table~V of the arXiv version
of~\cite{guan2022ganav}). Their groups (smooth, rough, bumpy, forbidden,
obstacle, background) differ from our six traversability groups, so these
figures show the gain from coarse grouping and are not directly comparable to
our group mIoU.

\subsection{Architectures and Pretraining}

U-Net~\cite{ronneberger2015unet} pairs a contracting encoder with a symmetric
expanding decoder and concatenates encoder features into the decoder at every
scale, which preserves fine spatial detail.
DeepLabV3+~\cite{chen2018deeplabv3plus} instead relies on atrous spatial
pyramid pooling (ASPP)~\cite{chen2017deeplabv3}, parallel dilated convolutions
that capture multi-scale context at reduced resolution, and adds a lightweight
decoder with a single low-level skip connection to refine boundaries. Both are
commonly paired with residual encoders~\cite{he2016resnet} pretrained on
ImageNet~\cite{deng2009imagenet}, which speeds convergence and helps when
labeled data is scarce. How well that transfer holds for low-viewpoint
natural terrain is one of the questions we test.

\subsection{Losses for Class Imbalance}

Weighted cross-entropy rescales each pixel's loss by a per-class weight. With
median-frequency balancing~\cite{eigen2015predicting}, the weight is the
median class frequency divided by the frequency of the pixel's class. Dice
loss~\cite{milletari2016vnet} directly optimizes a soft overlap measure that
is insensitive to how many pixels a class occupies. Focal
loss~\cite{lin2017focal} multiplies the cross-entropy term by
$(1-p_t)^\gamma$, which down-weights well-classified examples so that training
focuses on hard ones. Combining cross-entropy with Dice is a common way to
keep the stable gradients of the former while adding the class-balancing
effect of the latter. We compare all four under a fixed architecture.
